\documentclass[11pt]{article}

\usepackage[margin=1in]{geometry}
\usepackage{amsmath,amssymb,amsthm}
\usepackage[hidelinks]{hyperref}
\usepackage{listings}
\usepackage{booktabs}
\usepackage{xcolor}

\lstset{basicstyle=\ttfamily\small,breaklines=true,columns=fullflexible}
\setlength{\emergencystretch}{3em}

\newtheorem{theorem}{Theorem}
\newtheorem{lemma}[theorem]{Lemma}
\newtheorem{corollary}[theorem]{Corollary}
\newcommand{\NN}{\mathbb{N}}
\newcommand{\QQ}{\mathbb{Q}}
\newcommand{\bind}[2]{\binom{#1}{#2}}
\newcommand{\wedgep}[1]{\textstyle\bigwedge^{#1}}

\title{Two Proof Technologies for the Two-Families Theorem:\\
Permutations, Exterior Powers, and General Position in Lean 4}
\author{Louis Gestin\\\small Independent student, France\\\small \href{mailto:louis@gestin.net}{louis@gestin.net}}
\date{Independent technical report, July 2026}

\begin{document}
\maketitle

\begin{abstract}
We report a Lean~4 / Mathlib formalization of the Bollob\'{a}s--Lov\'{a}sz--Frankl
circle of two-families theorems, developed along two mathematically distinct proof
technologies. The first is a permutation-counting proof of Bollob\'{a}s's weighted
set-pairs inequality~\cite{bollobas1965} and its uniform corollary. The second is an
exterior-algebra proof of the Lov\'{a}sz--Frankl skew two-families theorem for
subspaces~\cite{lovasz1977}, bridged to the Frankl--Kalai skew theorem for finite
sets~\cite{frankl1982,kalai1984} by a moment-curve general-position construction
over~$\QQ$. All five principal theorems are fully machine-checked, each depending only
on the standard axioms \texttt{propext}, \texttt{Classical.choice}, \texttt{Quot.sound};
the development contains no \texttt{sorry}, \texttt{admit}, or custom axiom. The code
comprises roughly $1440$ lines across $13$ files ($51$ declarations, $37$ of them
theorems), builds cleanly in about $133$\,s when the pinned Mathlib cache is already available. Our thesis is that although the two proofs
use disjoint libraries, both produce the same abstract object---a triangular
certificate of independence---and we isolate several reusable components toward Mathlib:
a triangular linear-independence principle, vanishing/nonvanishing and concatenation
lemmas for decomposable exterior vectors, an exact division-free separation-event
count, and a moment-curve general-position theorem with span-intersection consequences.
The work is a formalization and introduces no new mathematics.
\end{abstract}

\section{Introduction}\label{sec:intro}

Two-families theorems are cornerstones of extremal combinatorics. In their simplest
uniform form they assert: if $A_1,\dots,A_m$ are $a$-sets and $B_1,\dots,B_m$ are
$b$-sets with $A_i\cap B_i=\varnothing$ and $A_i\cap B_j\neq\varnothing$ for $i\neq j$,
then $m\le\bind{a+b}{a}$. Bollob\'{a}s's weighted refinement~\cite{bollobas1965} drops
the uniformity and gives the exact rational bound
$\sum_i 1/\bind{|A_i|+|B_i|}{|A_i|}\le 1$. Lov\'{a}sz~\cite{lovasz1977} replaced sets by
subspaces and proved a \emph{skew} version (cross-intersection required only for
$i<j$), whose natural proof uses exterior algebra; Frankl~\cite{frankl1982} and
Kalai~\cite{kalai1984} transported the skew bound back to finite sets via a
general-position representation. Textbook expositions appear in
Babai--Frankl~\cite{babaifrankl1992}.

These proofs exhibit two genuinely different technologies. The weighted inequality is
proved by counting, over all linear orders of the ground set, the ``separation events''
in which each $A_i$ entirely precedes $B_i$; disjointness of these events across indices
yields the bound. The subspace theorem is proved by attaching to each pair a
decomposable exterior vector and a multiplication map, and invoking a triangular
linear-independence argument inside $\wedgep{a} V$, whose dimension is exactly
$\bind{a+b}{a}$. We formalize both, in Lean~4~\cite{lean4} on top of
Mathlib~\cite{mathlib2020}.

\paragraph{Thesis.} Our central observation is that the two proofs, despite using
disjoint parts of the library, converge on the same abstract certificate: a family of
objects that are ``nonzero on the diagonal and vanishing strictly below it.'' In the
counting proof this is a family of pairwise disjoint separation events of exactly
computed size; in the algebraic proof it is a family of exterior-multiplication test
maps feeding a single generic lemma
\lstinline{LinearIndependent.of_upperTriangular_maps}. Formalization makes the
correspondence precise rather than metaphorical.

\paragraph{Contributions.}
\begin{enumerate}
\item A complete, kernel-checked Lean development of the weighted and uniform Bollob\'as
  theorems, the Lov\'asz--Frankl subspace theorem, the moment-curve general-position
  theorem, and the Frankl--Kalai skew set theorem, along two independent proof routes
  (Sections~\ref{sec:count}--\ref{sec:bridge}).
\item Reusable infrastructure for which we did not locate a direct equivalent in Mathlib v4.28.0: triangular
  independence from test maps; exterior decomposable nonvanishing and concatenation; an
  exact, division-free count of order-separating enumerations; and moment-curve general
  position with span-intersection lemmas (Section~\ref{sec:reuse}).
\item An independent semantic audit, adversarial hypothesis testing, and a measured
  evaluation (Section~\ref{sec:eval}), supporting cautious novelty claims
  (Section~\ref{sec:related}).
\end{enumerate}

\section{Mathematical overview}\label{sec:overview}

Throughout, $\alpha$ is a finite ground type and $\iota$ a finite index type; $n=|\alpha|$.

\begin{theorem}[Weighted Bollob\'{a}s]\label{thm:weighted}
For $A_i,B_i\subseteq\alpha$ with $A_i\cap B_i=\varnothing$ for all $i$ and
$A_i\cap B_j\neq\varnothing$ for all $i\neq j$,
\[ \sum_{i} \frac{1}{\bind{|A_i|+|B_i|}{|A_i|}} \le 1. \]
\end{theorem}

\begin{corollary}[Uniform Bollob\'{a}s]\label{thm:uniform}
If additionally $|A_i|=a$ and $|B_i|=b$ for all $i$, then $|\iota|\le\bind{a+b}{a}$, and
the bound is attained by the complement family on $\{1,\dots,a+b\}$.
\end{corollary}

\begin{theorem}[Lov\'{a}sz--Frankl, subspaces]\label{thm:subspace}
Let $\dim V=a+b$ and $U_i,W_i\le V$ with $\dim U_i=a$, $\dim W_i=b$, $U_i\cap W_i=0$,
and $U_i\cap W_j\neq 0$ for $i<j$. Then $m\le\bind{a+b}{a}$.
\end{theorem}

\begin{theorem}[Frankl--Kalai, sets]\label{thm:fk}
Let $A_i,B_i\subseteq\alpha$ with $|A_i|=a$, $|B_i|=b$, $A_i\cap B_i=\varnothing$, and
$A_i\cap B_j\neq\varnothing$ for $i<j$. Then $m\le\bind{a+b}{a}$.
\end{theorem}

\noindent\emph{Short informal proofs.} For Theorem~\ref{thm:weighted}, a uniformly
random linear order of $\alpha$ separates $(A_i,B_i)$ with ``probability''
$1/\bind{|A_i|+|B_i|}{|A_i|}$; distinct separation events are disjoint (a crossing pair
forces contradictory orderings), so the reciprocals sum to at most $1$. We formalize
this by an \emph{exact integer count} rather than a probability. For
Theorem~\ref{thm:subspace}, decomposable vectors $u_i=\bigwedge U_i\in\wedgep{a}V$ and
multiplication by $w_j=\bigwedge W_j$ give maps with $u_i\wedge w_i\neq0$ (transversal
diagonal) and $u_i\wedge w_j=0$ for $i<j$ (a nonzero intersection makes the combined
family dependent); triangularity yields independence of the $u_i$, and
$m\le\dim\wedgep{a}V=\bind{a+b}{a}$. Theorem~\ref{thm:fk} follows by realizing sets as
moment-curve spans, which turns $\cap$-emptiness into $\oplus$ and $\cap$-nonemptiness
into a shared nonzero vector.

\section{Formalization architecture}\label{sec:arch}

The development is modular and splits into two routes that share only the predicate API
\lstinline{SetPairs.Basic}. Figure~\ref{fig:dep} shows the project-local import DAG
(generated from the \lstinline{import RequestProject.*} lines; verified acyclic in
\texttt{DEPENDENCY\_AUDIT.md}).

\begin{figure}[t]
\centering
\begin{verbatim}
                         SetPairs/Basic
                         /            \
          PermutationOrders          SkewSetPairs
                |                      /        \
          SeparationEvents           /          \
                |             MomentCurve   TwoFamiliesSubspaces
          WeightedBollobas                 /     |         \
                |                Triangular  ExteriorDecomp  ExteriorMult
             Uniform                        Independence
                \                              /
                 \____________ Main _________/    (+ Examples)
\end{verbatim}
\caption{Project module dependency graph. Left branch: permutation counting. Right
branch: exterior powers and general position. The two branches meet only at
\texttt{SetPairs/Basic} (predicates) and at \texttt{Main}.}
\label{fig:dep}
\end{figure}

\begin{itemize}
\item \texttt{SetPairs/}: \texttt{Basic} (predicates), \texttt{PermutationOrders},
  \texttt{SeparationEvents}, \texttt{WeightedBollobas}, \texttt{Uniform}.
\item \texttt{LinearAlgebra/}: \texttt{TriangularIndependence},
  \texttt{ExteriorDecomposable}, \texttt{ExteriorMultiplication},
  \texttt{TwoFamiliesSubspaces}, \texttt{MomentCurve}.
\item \texttt{SkewSetPairs}, \texttt{Examples}, \texttt{Main}.
\end{itemize}

\section{Exact permutation counting}\label{sec:count}

An \emph{order enumeration} of $\alpha$ is an equivalence $e:\alpha\simeq\mathrm{Fin}\,n$;
there are $n!$ of them (\lstinline{card_orderEnumerations}, from
\lstinline{Fintype.card_equiv}). The pair $(A,B)$ is \emph{separated} by $e$ when every
element of $A$ precedes every element of $B$:
\begin{verbatim}
def Separates (e : OrderEnumeration a) (A B : Finset a) : Prop :=
  forall {x}, x in A -> forall {y}, y in B -> e x < e y
\end{verbatim}
and $\mathrm{separationEvent}(A,B)$ is the finset of separating enumerations.

\paragraph{The exact count.} The heart of the argument is
\[ |\mathrm{separationEvent}(A,B)|\cdot\bind{|A|+|B|}{|A|} = n!, \]
which we obtain from the closed form (\lstinline{card_separationEvent_eq})
\[ |\mathrm{separationEvent}(A,B)| = \bind{n}{a+b}\,a!\,b!\,(n-(a{+}b))!, \]
with $a=|A|$, $b=|B|$, $S=A\cup B$. This count is proved by a genuine
\emph{ordering-factorization} equivalence: a separating enumeration is determined by
three independent pieces:
\begin{itemize}
\item the $(a{+}b)$-element position set $P=e(S)\subseteq\mathrm{Fin}\,n$ (contributing
  $\bind{n}{a+b}$);
\item the relative order of $S$ inside $P$, in which $A$ precedes $B$ (contributing
  $a!\,b!$);
\item the relative order of $S^{\mathsf c}$ inside $P^{\mathsf c}$ (contributing
  $(n-a-b)!$).
\end{itemize}
Concretely:
\begin{itemize}
\item \lstinline{card_separates_full}: when $A\cup B=\mathrm{univ}$, the separating
  enumerations biject with $(A\simeq\mathrm{Fin}\,a)\times(B\simeq\mathrm{Fin}\,b)$,
  giving $a!\,b!$. The bijection is built from two explicit injections whose
  well-definedness rests on the order-statistics facts that separation forces every
  $A$-element below rank $a$ (\lstinline{separates_val_lt}) and every $B$-element at
  rank $\ge a$ (\lstinline{card_le_separates_val}).
\item \lstinline{card_separates_toP}: transports this count along the order isomorphism
  $P\simeq_o\mathrm{Fin}(a{+}b)$ to an arbitrary ordered position set $P$.
\item \lstinline{fiber_card_eq}: for a fixed $P$, the enumerations with $e(S)=P$ biject
  with (relative orders of $S$)~$\times~(S^{\mathsf c}\simeq P^{\mathsf c})$.
\item \lstinline{card_separationEvent_eq}: sums the fibre counts $a!\,b!\,(n-(a{+}b))!$
  over the $\bind{n}{a+b}$ position sets, via
  \lstinline{Finset.card_eq_sum_card_fiberwise}.
\end{itemize}
The factorial identity then follows \emph{without any natural-number division}, from
$\bind{a+b}{a}a!b!=(a+b)!$ and $\bind{n}{a+b}(a+b)!(n-(a{+}b))!=n!$
(\lstinline{Nat.choose_mul_factorial_mul_factorial}, applied twice).

\paragraph{Why not probability.} A direct probabilistic phrasing would introduce
$\mathbb{R}$-valued measures and division, whose formal handling is heavier than the
integer identity above and obscures the fibre structure that is the real proof. The
integer count is also what the regression tests can check exactly
(Section~\ref{sec:eval}).

\paragraph{Disjointness and assembly.} Separation events are pairwise disjoint: if
$x\in A_i\cap B_j$ and $y\in A_j\cap B_i$, then separating $(A_i,B_i)$ forces
$e(x)<e(y)$ while separating $(A_j,B_j)$ forces $e(y)<e(x)$
(\lstinline{separationEvents_disjoint}). This step is exactly where cross-intersection
in \emph{both} directions is used. Summing the disjoint cardinalities gives
$\sum_i|E_i|\le n!$, and since $1/\bind{|A_i|+|B_i|}{|A_i|}=|E_i|/n!$ the weighted
inequality (\lstinline{weighted_bollobas}) follows. The uniform corollary
(\lstinline{uniform_bollobas}) factors out the constant reciprocal and divides;
sharpness is the complement family on $\mathrm{Fin}(a+b)$
(\lstinline{uniform_bollobas_sharp}).

\section{Triangular independence and exterior products}\label{sec:ext}

\paragraph{Triangular independence.} If linear maps $T_j$ satisfy $T_j(u_j)\neq0$ and
$T_j(u_i)=0$ for $i<j$, then the $u_i$ are linearly independent
(\lstinline{LinearIndependent.of_upperTriangular_maps}):
\begin{verbatim}
theorem LinearIndependent.of_upperTriangular_maps
    (u : Fin m -> M) (T : Fin m -> M ->L[K] N)
    (hdiag : forall j, T j (u j) != 0)
    (hzero : forall {i j}, i < j -> T j (u i) = 0) :
    LinearIndependent K u
\end{verbatim}
The proof takes a vanishing combination $\sum_i g_i u_i=0$, lets $j$ be the
\emph{largest} index with $g_j\neq0$, and applies $T_j$: terms with $i<j$ die by
$\mathtt{hzero}$, terms with $i>j$ die by maximality, leaving $g_j\,T_j(u_j)=0$, a
contradiction since $K$ is a field. The maximal-index choice is matched to the
upper-triangular orientation; a dual \lstinline{of_lowerTriangular_maps} uses the
minimal index.

\paragraph{Decomposable vectors.} The alternating product $\iota\mathrm{Multi}\,K\,r\,v$
vanishes exactly when $v$ is linearly dependent: the vanishing direction is
\lstinline{AlternatingMap.map_linearDependent}, and the nonvanishing direction
(\texttt{\(\iota\)Multi\_ne\_zero\_of\_linearIndependent}, both in the exterior power and the
exterior algebra) is obtained from Mathlib's
\texttt{\(\iota\)Multi\_family\_linearIndependent\_field} applied to the unique full subset.
Concatenation
$\iota\mathrm{Multi}\,a\,u\cdot\iota\mathrm{Multi}\,b\,w=\iota\mathrm{Multi}\,(a{+}b)\,(u\mathbin{+\!\!+}w)$
(\texttt{\(\iota\)Multi\_mul\_\(\iota\)Multi}) turns products of decomposables into the decomposable of
the appended family; its proof is an element-wise
\lstinline{List.ofFn}/\lstinline{Fin.append} identity.

\paragraph{Subspace theorem.} For each $i$ we pick bases of $U_i$ and $W_i$ (via
\lstinline{Module.finBasisOfFinrankEq}, using $\dim U_i=a$, $\dim W_i=b$), forming
$u_i\in\wedgep{a}V$ and a test map $T_j=(\,\cdot\,*w_j)\circ\iota$ where $w_j$ is the
decomposable of $W_j$. Then $T_j(u_i)$ is the decomposable of the concatenated bases,
nonzero on the diagonal (since $U_i\cap W_i=0$ makes the union independent,
\lstinline{append_wedge_ne_zero}) and zero for $i<j$ (since $U_i\cap W_j\neq0$ makes it
dependent, \lstinline{append_wedge_eq_zero}). Triangular independence bounds
$m\le\dim\wedgep{a}V=\bind{a+b}{a}$ via \lstinline{exteriorPower.finrank_eq}. The
degenerate cases $a=0$, $b=0$, $m=0$ require no special handling: empty bases and
$\iota\mathrm{Multi}\,K\,0=1$ are absorbed uniformly---a case split that appeared
necessary on paper turned out to be an artifact of informal reasoning.

\section{Moment-curve bridge}\label{sec:bridge}

The moment-curve vector is $v_x=(1,t_x,\dots,t_x^{d-1})\in\QQ^d$ for an injective label
$t:\alpha\hookrightarrow\QQ$. Any at most $d$ distinct such vectors are linearly
independent (\lstinline{momentCurve_linearIndependent}); we prove this by polynomial
interpolation---the polynomial $\prod_{j\neq i}(X-t_j)$, of degree $|s|-1<d$, isolates
index $i$---rather than by expanding the Vandermonde determinant, which is cleaner in
Lean and, crucially, handles the \emph{rectangular} case $|s|<d$ directly without any
square-matrix reduction. As consequences we obtain the span dimension
$\dim\operatorname{span}\{v_x:x\in S\}=|S|$ for $|S|\le d$
(\lstinline{finrank_spanMomentCurve}), triviality of
$\operatorname{span}(S)\cap\operatorname{span}(T)$ for disjoint $S,T$ with $|S|+|T|\le d$
(\lstinline{inf_spanMomentCurve_eq_bot}), and nontriviality when $S\cap
T\neq\varnothing$ (\lstinline{inf_spanMomentCurve_ne_bot}, using that each $v_x$ is
nonzero because its zeroth coordinate is $1$). The Frankl--Kalai theorem then labels the
ground type injectively into $\QQ$ (via \lstinline{Fintype.equivFin} composed with
$\NN\hookrightarrow\QQ$), realizes each set as a moment-curve span in $\QQ^{a+b}$, and
applies the subspace theorem. The ground set may be strictly larger than $a+b$; only the
ambient dimension $d=a+b$ is constrained.

\section{Reusable infrastructure and Mathlib gaps}\label{sec:reuse}

Table~\ref{tab:correspondence} maps the central mathematical objects to their Lean
representations. Several components are generic and were not located in Mathlib v4.28.0 during our documented search
(\texttt{MATHLIB\_PR\_PLAN.md}): triangular independence from test maps; exterior
decomposable nonvanishing and concatenation (Mathlib has only the vanishing direction);
the exact separation-event count; and moment-curve general position with
span-intersection lemmas. Mathlib does supply the exterior-power API and its dimension
$\dim\wedgep{a}V=\bind{\dim V}{a}$ (\lstinline{exteriorPower.finrank_eq}), the
Vandermonde determinant, and the LYM inequality (a different chain-counting technique).

\begin{table}[t]
\centering\small
\begin{tabular}{@{}llll@{}}
\toprule
Mathematical object & Lean representation & Key declaration & Main difficulty \\
\midrule
Total ordering & $\alpha\simeq\mathrm{Fin}\,n$ & \texttt{OrderEnumeration} & counting bijections \\
Separation event & filtered \texttt{Finset} & \texttt{separationEvent} & exact fibre count \\
Exterior decomposable & $\iota$\texttt{Multi}$\,K\,r\,v$ & \texttt{\(\iota\)Multi\_ne\_zero\_*} & nonvanishing direction \\
Exterior mult. test & \texttt{mulRight} $\circ$ subtype & \texttt{\(\iota\)Multi\_mul\_\(\iota\)Multi} & append identity \\
Moment-curve vector & $\mathrm{Fin}\,d\to\QQ$ & \texttt{momentCurve} & rectangular independence \\
Set span & \texttt{Submodule.span} & \texttt{spanMomentCurve} & $\cap$ vs.\ $\oplus$ \\
Skew cross-intersection & $i<j\to(A_i\cap B_j)\neq\varnothing$ & \texttt{SkewCrossIntersecting} & orientation \\
\bottomrule
\end{tabular}
\caption{Informal-to-Lean correspondence (Table A).}
\label{tab:correspondence}
\end{table}

\begin{table}[t]
\centering\small
\begin{tabular}{@{}lllccl@{}}
\toprule
Theorem & Method & Lean declaration & Checked & Sharp & Reusable deps \\
\midrule
Weighted & counting & \texttt{weighted\_bollobas} & yes & --- & sep.\ count \\
Uniform & counting & \texttt{uniform\_bollobas} & yes & yes & weighted \\
Triangular & lin.\ alg. & \texttt{of\_upperTriangular\_maps} & yes & --- & (generic) \\
Subspace & exterior & \texttt{lovasz\_frankl\_subspaces} & yes & --- & triangular, exterior \\
Moment curve & interp. & \texttt{momentCurve\_linearIndependent} & yes & --- & (generic) \\
Frankl--Kalai & bridge & \texttt{frankl\_kalai\_skew} & yes & yes & subspace, moment \\
\bottomrule
\end{tabular}
\caption{Theorem coverage (Table B). ``Sharp'' = a matching extremal family is
formalized; the uniform/skew bound is attained by the complement construction.}
\label{tab:coverage}
\end{table}

\section{Evaluation}\label{sec:eval}

\paragraph{Versions.} Lean \texttt{v4.28.0}; Mathlib pinned at \texttt{v4.28.0} (commit
\texttt{8f9d9cff}). Project clean build with the pinned Mathlib artifacts already cached:
\textbf{success, 8039 jobs, $\approx 133$\,s} wall-clock in our environment. A cold reconstruction first downloads the pinned dependencies and obtains the Mathlib cache, so it takes substantially longer.

\paragraph{Size.} $13$ Lean files, $1438$ total lines ($1175$ non-blank), $51$
declarations: $37$ theorems/lemmas, $10$ definitions, $2$ abbreviations, $2$ instances.
By module the largest is \texttt{SeparationEvents.lean} ($481$ lines), followed by
\texttt{MomentCurve.lean} ($179$) and \texttt{TwoFamiliesSubspaces.lean} ($144$).

\paragraph{Axioms.} Every principal theorem's \lstinline{#print axioms} lists exactly
\lstinline{[propext, Classical.choice, Quot.sound]}; no \lstinline{sorryAx}. The
repository contains no \lstinline{sorry}, \lstinline{admit}, custom \lstinline{axiom},
\lstinline{unsafe}, or \lstinline{@[implemented_by]} (verified by
\texttt{scripts/check\_forbidden.sh}).

\paragraph{Regression tests.} \texttt{scripts/verify\_small\_cases.py} exhaustively
checks, in exact rational arithmetic: the weighted inequality on ground sets of size
$2$--$5$; the uniform bound for admissible $(a,b)$ on ground sets of size $\le 6$; the
skew bound on ground sets of size $\le 4$; and the sharpness construction
(\texttt{OVERALL: PASS}). \texttt{scripts/adversarial\_checks.py} confirms with exact
counterexamples that each principal hypothesis is load-bearing
(\texttt{ADVERSARIAL\_TESTS.md}). Small Lean \lstinline{decide}/\lstinline{norm_num}
checks (\texttt{Examples.lean}) corroborate specific binomial values, the closed-form
count $\bind{n}{a+b}a!b!(n-(a{+}b))!$ over five boundary cases (empty $A$/$B$, elements
outside $A\cup B$, and $A\cup B=\mathrm{univ}$), and two adversarial sums exceeding $1$.
The computational checks are regression tests, not substitutes for the proofs.

\paragraph{Reusable fraction.} The generic infrastructure
(Sections~\ref{sec:ext}--\ref{sec:reuse}: triangular independence, exterior
nonvanishing/concatenation, moment-curve general position, separation-event counting)
is application-independent and constitutes the bulk of the extraction candidates in
\texttt{MATHLIB\_PR\_PLAN.md}.

\paragraph{Failed and revised approaches.} Moment-curve independence was first attempted
through the rectangular Vandermonde matrix, then replaced by the cleaner
polynomial-interpolation argument. The subspace theorem was initially planned with
explicit case splits on $a=0$, $b=0$, $m=0$; these proved unnecessary once the argument
was routed through the exterior algebra directly. The separation-event count was
initially phrased with natural-number division and re-derived as an exact product
identity to avoid divisibility side-conditions.

\section{Related work}\label{sec:related}

The mathematical originals are Bollob\'as~\cite{bollobas1965} (weighted/uniform
set-pairs), Lov\'asz~\cite{lovasz1977} (exterior-algebra subspace method), and
Frankl~\cite{frankl1982}/Kalai~\cite{kalai1984} (skew set extensions), with the
exposition of Babai--Frankl~\cite{babaifrankl1992}. Mathlib~\cite{mathlib2020} provides
the exterior-power and graded-exterior-algebra API, the Vandermonde determinant, and the
LYM inequality, but---to the best of our documented search
(\texttt{NOVELTY\_AUDIT.md})---no Bollob\'as set-pairs inequality, uniform/skew
two-families theorem, or Lov\'asz subspace theorem. The local Mathlib source was searched directly, while external indexes and libraries were surveyed to the extent available in the audit environment; a live re-check remains advisable before making a stronger novelty claim. We did not locate a prior formalization combining the weighted permutation proof, the Lov\'asz--Frankl subspace theorem, and the
moment-curve reduction to the Frankl--Kalai skew theorem in one development, and make no
stronger claim (in particular not ``the first formalization'' of any individual
theorem, since live external search was unavailable).

\section{Lessons and limitations}\label{sec:lim}

\paragraph{Lessons.} (i) Exact integer identities beat probabilistic phrasing in a proof
assistant, both for the proof skeleton and for exact regression testing. (ii) Setting up
the algebra uniformly can make degenerate cases dissolve rather than require case
splits. (iii) Rectangular general position must be genuinely proved (via interpolation),
not assumed by analogy with the square Vandermonde matrix. (iv) Subtype/cardinality
transport (between $\uparrow S$, $\mathrm{Fin}\,|S|$, and position sets) is a recurring
cost centre in the counting proof, handled by explicit order isomorphisms.

\paragraph{Limitations.} The development relies on classical choice and finite ground
and index types. Equality/extremal classification beyond the single complement
construction is not formalized. The moment-curve layer is stated over $\QQ$ rather than
an arbitrary (infinite) field. As the entire project workflow was substantially AI-assisted---including planning, proof development, code, testing, writing, translation, auditing, and repository preparation---kernel-checking
guarantees correctness only \emph{relative to the stated theorems}; semantic fidelity is
argued separately (\texttt{INDEPENDENT\_PROOF\_AUDIT.md}, \texttt{SEMANTIC\_AUDIT.md})
and every claim remains the responsibility of the human author
(\texttt{AI\_USE\_STATEMENT.md}).

\section{Conclusion}\label{sec:concl}

We have formalized, along two distinct proof technologies, the
Bollob\'{a}s--Lov\'{a}sz--Frankl two-families circle in Lean~4---the weighted and
uniform Bollob\'{a}s theorems by permutation counting, and the Lov\'{a}sz--Frankl
subspace and Frankl--Kalai skew set theorems by exterior powers and moment-curve general
position---together with reusable algebraic and combinatorial infrastructure. All five
principal theorems are kernel-checked under the standard axioms. Beyond the specific
theorems, the development isolates a common triangular-certificate structure underlying
two superficially unrelated proofs, and packages several components in Mathlib-ready
form. The work claims no new mathematics.

\section*{Artifact availability}
The complete Lean sources, scripts, audit documents, LaTeX sources, and both language
versions of this report are distributed in the accompanying project archive. The development
is pinned to Lean \texttt{v4.28.0} and Mathlib \texttt{v4.28.0}; the principal build command is
\lstinline{lake build}.

\bibliographystyle{plain}
\bibliography{references}

\end{document}
